\documentclass[11pt,a4paper]{article}
\usepackage{fontspec}
\setmainfont{DejaVu Serif}
\usepackage{amsmath,amssymb,graphicx,booktabs,geometry,hyperref}
\geometry{margin=22mm}
\sloppy
\setlength{\emergencystretch}{3em}
\graphicspath{{results/}}
\begin{document}
\title{Що з’ясовано та що дослідити далі}\author{}\date{27 жовтня 2026}\maketitle

Формулу відновлено з авторських рівнянь:
$\phi_{n,r}(s)=(1-s^2/r^2)_+^{n+1}$, гладкість --- $C^n$.
Точне проходження через вузли та повний ранг не гарантують
точності між вузлами: для константи при $h=0.25$, $n=2$,
$r/h=2$ і запасі $r$ вибіркова максимальна похибка дорівнює $0.0744$.

Побудовано чотири точні 2D-еталони та сферу й тор у 3D.
Кути, медіальні множини та критичні офсети треба розглядати окремо.
Дуже тонкий офсет тора розпадається навіть в еталонному сітковому
поданні: це обмеження роздільності, а не доведена помилка моделі.

У серії повернутого квадрата RMSE смуги при
$r/h=1.5,2,3$ становить відповідно $0.0126$, $0.0340$, $0.1352$.
При $r/h=3$ виявлено 6 сіткових компонент моделі замість 1.
За сталого $r=0.5$ ущільнення з $h=0.32$ до $0.16$ збільшило
RMSE приблизно з $0.0184$ до $0.0684$. Отже, радіус і густину
потрібно досліджувати разом. Це висновки конкретних конфігурацій.

\textbf{Пункт 4 залишається актуальним:} відділити форми для
налаштування від незалежної фінальної перевірки. Поточні приклади
є діагностичними. Потрібні кілька seeds і взаємодії факторів.

\textbf{Пункт 5 виконано частково:} є похибки поля, знака,
компоненти, ранг малих матриць і залишки. Далі потрібні прямі
відстані між офсетними поверхнями, градієнти поза медіальними
множинами, топологія отворів та перевірка кількох роздільностей.
RMSE поля не є відстанню Гаусдорфа між поверхнями.

\textbf{Пункт 6 актуальний:} побудувати карту відмов і криві
Парето за пам’яттю, точністю та часом. Для цього потрібні
вимірювання запитів і пікової пам’яті та порівняння з Wendland
$C^2/C^4$, білінійною/трилінійною й кубічною інтерполяцією.

Перший наступний крок --- похибка самих офсетів при фіксованих
пам’яті й бюджеті налаштування. Окремо можна перевірити доповнення,
яке точно відтворює константу й афінні функції. Це вже модифікація
початкового методу. Нинішні результати мають один seed і початкову
роздільність $35^3$ для 3D; доведення збіжності та повної перевірки
літературної новизни поки немає.

\end{document}
